
\newtoggle{withbonus} \settoggle{withbonus}{false} % set true to show bonus points in grade table

% setting underlying course name (Grundlagenfach Mathematik, §-Unterricht, \ldots)
\newcommand{\mycourse}{Informatik}
% name of the class
\newcommand{\myclass}{\emoji{clapping-hands}\faPeace\emoji{handshake}~Klasse 2g~\emoji{handshake}\faPeace\emoji{clapping-hands}}
% set date
\newcommand{\mydate}{Donnerstag, 12. Juni 2025}

\title{\textbf{Randomisierte Algorithmen, Kommunikationsprotokoll}}
\author{Thomas Graf\\
	\mycourse, \myclass}
\date{\mydate}
\maketitle
\thispagestyle{headandfoot}

\begin{center}
	\fbox{\parbox{140mm}{\centering
			\begin{itemize}
				\item Dauer der Prüfung: $50$ Minuten
				      %\item \textbf{Zeige in jeder Aufgabe unbedingt Deine Schritte!}
				\item erlaubte Hilfsmittel:
				      \begin{itemize}
					      \item Schreibzeug
				      \end{itemize}
			\end{itemize}}}
\end{center}
\vspace{10mm}

\begin{center}
	Vorname: \rule{45mm}{0.8pt}\qquad Nachname: \rule{45mm}{0.8pt}
\end{center}
\vspace{20mm}

\begin{center}
	\iftoggle{withbonus} {
		\combinedgradetable[h][questions]
	} {
		\gradetable[h][questions]
	}
\end{center}
\vspace{15mm}
\begin{center}
	Note:
	\vspace{10mm}

	\rule{8mm}{1.2pt}
\end{center}
%\thispagestyle{empty}
\clearpage


\begin{questions}
	% 2g
	\question
	\begin{parts}
		\part[1] Bestimmen Sie $\text{Nummer}(0101101)$.
		\part[1] Bestimmen Sie binäre Wort $x$, sodass $\text{Nummer}(x) = 82$.
		\part[1] Berechnen Sie $\abs{\text{Bin}(3099)}$.
		\part
		Sei $y$ die binäre Darstellung ($2$-adische Darstellung) der drittgrössten Zahl, welche mit $9$ Bits dargestellt werden kann.
		\begin{subparts}
			\subpart[1] Wie sieht $y$ aus?
			\subpart[1] Bestimmen Sie $\text{Nummer}\lr{y}$.
		\end{subparts}
		\part[1] Berechnen Sie (stellen Sie ohne Logarithmus dar) die Zahl $\ceil{\log_5\lr{89}}$.
		\part[1] Donald Knuth möchte unser randomisierte Kommunikationsprotokoll für die beiden Datensätze $x := 001110$ und $y := 1110$ durchführen.
		Was ist hier das Problem? Welchen Rat würden Sie Donald geben? Argumentieren Sie präzise.
		\part[1] Bestimmen Sie $\text{Prim(37)}$.
		\part[1] Nennen Sie (konkrete) positive natürliche Zahlen $x, y$ und $z$, sodass zwar die Relation
		\begin{align*}
			x < y < z
		\end{align*}
		gilt, aber nicht die Relation
		\begin{align*}
			\abs{\text{Bin}(x)} < \abs{\text{Bin}(y)} < \abs{\text{Bin}(z)}.
		\end{align*}
	\end{parts}
	\droptotalpoints


	\question
	\begin{parts}
		\part[2]
		Gegeben sind die beiden Datensätze
		\begin{itemize}
			\item $x := 001001$ in Zürich und
			\item $y := 101010$ in Boston.
		\end{itemize}
		Für welche Primzahlen $p \leq n^2$ würde das Kommunikationsprotokoll für diese beiden Datensätze $x$ und $y$ falsch entscheiden? Begründen Sie Ihre Antwort.
		\part[1] Es seien nun $x := 101110$ und $y := 101111$. Kann es sein, dass das Kommunikationsprotokoll für diese Datensätze eine falsche Entscheidung treffen wird? Begründen Sie Ihre Antwort.
	\end{parts}
	\droptotalpoints


	\question Gegeben seien zwei Datensätze $x$ und $y$ mit einer jeweiligen Länge von $10^{20}$ Bits.
	\begin{parts}
		\part[2] Geben Sie eine vernünftige (möglichst kleine) obere Schranke für den Kommunikationsaufwand des Kommunikationsprotokolls zum Abgleich von $x$ und $y$ an.
		\part[1] Wie viele Bits müsste ein deterministischer Algorithmus zur Beantwortung der Frage $x\stackrel{?}{=}y$ mindestens kommunizieren?
	\end{parts}
	\droptotalpoints

	\clearpage
	\question[2]
	Die Analyse der Fehlerwahrscheinlichkeit $r$ des Protokolls wird die Abschätzung
	\begin{align*}
		r \leq \frac{n-1}{\text{Prim}(n^2)}\approx\frac{n-1}{n^2/\ln(n^2)}\leq\frac{\ln(n^2)}{n}
	\end{align*}
	erfüllen. Dabei ist $n$ die Länge der jeweiligen Datensätze (in Bits). Erklären Sie, warum wir im Protokoll eine Primzahl $p \leq n^2$ wählen und nicht lediglich eine Primzahl $p \leq n$.

	Tipp: Wie verhält sich die Fehlerwahrscheinlichkeit für wachsendes $n$?

	\question[2]
	Gegeben ist die (korrekte und vollständige) Ihnen vertraute Python-Funktion \pythoninline{sieb}:
	\begin{lstlisting}[language=Python,caption=Sieb,numbers=none]
import math

def sieb(n):
    primzahlen = []
    nicht_gestrichen = (n + 1) * [True]

    i = 2
    # abgerundete Quadratwurzel von n
    wurzel_n = int(math.sqrt(n))
    while i <= wurzel_n:
        if nicht_gestrichen[i]:
            j = i * i
            while j <= n:
                nicht_gestrichen[j] = False
                j += i
        i += 1

    # sammle alle nicht gestrichenen Zahlen
    i = 2
    while i <= n:
        if nicht_gestrichen[i]:
            primzahlen.append(i)
        i += 1
    return primzahlen
\end{lstlisting}
	Verwenden Sie die Funktion \pythoninline{sieb} als Hilfsfunktion um eine Python-Funktion \pythoninline{prime_gaps(n)} zu definieren, welche eine Liste der $(n - 1)$ Primzahllücken (prime gaps) zwischen den ersten $n$ Primzahlen printet. Vervollständigen Sie die gegebene \textbf{Vorlage}.

	\begin{lstlisting}[language=Python,caption=Beispiele,numbers=none]
prime_gaps(20) # printet die Liste: [1, 2, 2, 4, 2, 4, 2]
prime_gaps(7)  # printet die Liste: [1, 2, 2]
\end{lstlisting}

	\clearpage

	\begin{lstlisting}[language=Python,caption=prime gaps (noch zu vervollständigen),numbers=none]
def prime_gaps(n):
    # Get all primes up to n
    primes = sieb(n)

    # Calculate the gaps between consecutive primes
    gaps = []

    ..................................... 
    

    ..................................... 
    

    ..................................... 


    ..................................... 
    

    ..................................... 
    

    print(gaps)
\end{lstlisting}

	% \bonusquestion[2]
	% Es sei $n$ eine beliebige ganze Zahl. Betrachten Sie die endliche Folge ganzer Zahlen
	% \begin{align*}
	% 	(n+1)!+2,(n+1)!+3,\ldots, (n+1)!+(n+1)
	% \end{align*}
	% Erklären Sie, warum
\end{questions}
